\section{Controllers (ACC, CACC)}
\label{sec:sota_controller}

Longitudinal control is at the core of platooning: it determines inter-vehicle distance, response speed to disturbances, and ultimately the safety of the system. A first fundamental distinction is between controllers based only on local sensors (ACC) and cooperative controllers (CACC), which exploit measurements and state information shared through V2V communication to improve responsiveness and performance \cite{ploeg2011design,oncu2014cooperative,segata2016safe}. This distinction is central because the performance gain of platooning derives precisely from the possibility of reducing safety margins with respect to a traditional ACC, but that reduction is possible only if channel availability and quality remain compatible with the controller requirements.

In the case of ACC, the vehicle regulates speed and distance with respect to its predecessor, typically using radar and a \emph{constant time-headway} policy. This policy increases the safety distance with speed and is robust, but it limits the benefits in terms of traffic density and energy savings. Segata's dissertation \cite{segata2016safe} also highlights the role of \emph{actuation lag}, namely the delay between a command and the actual response of the vehicle, which directly constrains the controller parameters and the minimum maintainable distance under stable conditions.

For CACC, the literature proposes different control topologies and spacing policies. A useful classification axis distinguishes between \emph{predecessor-following} controllers, which primarily exploit the immediately preceding vehicle, and \emph{leader + predecessor-following} controllers, which also integrate information from the platoon leader \cite{ploeg2011design,segata2016safe}. The former are generally closer to the logic of an advanced ACC and lend themselves to gradual transitions toward non-cooperative modes; the latter can offer better performance and, in some cases, support \emph{constant spacing} policies, but they require richer and more reliable connectivity.

The main theoretical objective is \emph{string stability}, that is, the ability of the platoon not to amplify disturbances in position, speed, or acceleration toward the tail of the platoon \cite{besselink2017string,segata2016safe}. This property depends jointly on controller parameters, the vehicle dynamic model, the spacing policy, and communication quality in terms of latency, jitter, and losses. In practical terms, a controller that is nominally valid may lose its desired properties if it is informed too infrequently; for this reason, analyzing the controller alone, without a realistic channel model, is often insufficient \cite{giordano2019joint}.

Among the controllers most relevant to the rest of the thesis are the Ploeg CACC and the PATH CACC, both also used in the multi-technology work of Segata \emph{et al.}\cite{segata2023multi-technology}. The Ploeg controller is often employed as an intermediate stage during degradation because it retains logic close to an advanced ACC, with predecessor information received via radio, and supports constant time-headway policies; the PATH controller, by contrast, is more demanding in terms of communication but allows more compact and higher-performance configurations, and is therefore appropriate when connectivity is more reliable \cite{ploeg2011design,segata2016safe,segata2023multi-technology}.

This is where the need for \emph{graceful degradation} mechanisms arises, in which the system transitions between different controllers, and different gap parameters, as a function of the observed degradation. The literature shows that direct transitions to ACC can compromise comfort and safety, whereas controlled degradation strategies and joint network/control design make it possible to preserve safety constraints even in the presence of persistent or temporary losses \cite{ploeg2015graceful,giordano2019joint,segata2023multi-technology}.
